\begin{abstract}
Decentralized storage networks (DSNs) are storage systems powered by permissionless nodes.
Data placement in DSNs must tolerate not only storage-device failures but also adversarial behavior that targets data availability.
Byzantine nodes introduce unique challenges due to collusion and adaptive attacks.
They can target specific data blocks by clustering within a block's placement group, reducing the number of rational nodes and weakening failure tolerance.
In this work, we propose a global defense against Byzantine nodes across all placement groups.
We introduce a node-centric approach that guarantees stable incentives for rational nodes regardless of the number of Byzantine nodes in their placement groups.
Building on this approach, we design \sys, a DSN that uses sampling-based data placement with verifiable randomness.
Compared with prior DSNs, this placement strategy allows \sys to scale simultaneously in storage volume, on-chain footprint, and Byzantine tolerance.
Our preliminary results show that \sys achieves the desired availability with scalable storage overhead while maintaining scalable fault tolerance.
\end{abstract}